\chapter{磁冻结、X点电场与能量：重联究竟改变了什么}
\label{hall:ch:freezing}
\section{不要把磁力线当作可以剪断的物质细绳}
磁力线是某个时刻满足$\dd\vct x/\dd l\parallel\B$的几何曲线。它由磁场定义，没有给每条线附上唯一粒子编号。所谓“磁冻结”，是一个关于随流曲面磁通保持的精确定理；在足够光滑、无奇点且边界适当时，它允许我们将磁结构随某个流速输运。

磁重联最有用的描述是磁连接关系与磁通转移发生变化，并伴随场能向粒子能和输运通道转化。示意图中“切断再接上”的线条有助于辨认连接变化，却没有告诉我们欧姆定律中究竟由哪一项支持非理想电场。下面从Faraday定律出发建立这个判断。

\section{移动曲面的磁通导数：固定面与扫过的窄条}
设曲面$S(t)$的边界$C(t)$以速度$\vct v_f$移动，法向与边界方向满足右手规则。磁通定义为
\begin{equation}
\Phi_B(t)=\int_{S(t)}\B(\vct x,t)\cdot\dd\vct S.
\end{equation}
在小时间$\dd t$内，磁通变化有两部分。

第一部分是在原曲面上磁场本身变化：
\begin{equation}
\delta\Phi_{\rm field}=\dd t\int_{S(t)}\partial_t\B\cdot\dd\vct S
=-\dd t\oint_{C(t)}\E\cdot\dd\vct l.
\end{equation}
这里用了Faraday定律和Stokes定理。第二部分来自边界移动扫过的窄条面积。以定向线元$\dd\vct l$表示边界，这个窄条的定向面积为$\vct v_f\dd t\times\dd\vct l$，因此
\begin{align}
\delta\Phi_{\rm motion}
&=\dd t\oint_C\B\cdot(\vct v_f\times\dd\vct l)\\
&=-\dd t\oint_C(\vct v_f\times\B)\cdot\dd\vct l.
\end{align}
二者相加，除以$\dd t$再取极限：
\begin{equation}
\boxed{\frac{\dd\Phi_B}{\dd t}
=-\oint_{C(t)}(\E+\vct v_f\times\B)\cdot\dd\vct l.}
\label{hall:eq:moving-flux}
\end{equation}
检查符号的方法是取均匀$\B=B\hat{\vct z}$、一个向外扩大的圆环。$\vct v_f$径向向外，$\vct v_f\times\B$沿顺时针；若正边界是逆时针，右边应为正，确实给出面积增大时磁通增大。

\subsection{冻结的条件是闭合积分，而不只是某一点电场为零}
若
\begin{equation}
\E+\vct v_f\times\B=-\nabla\Phi,
\end{equation}
且势函数在相关区域单值，则闭合回路积分为零，磁通冻结。一个常用充分条件是$\nabla\times(\E+\vct v_f\times\B)=0$，再配合单连通区域等全局条件。理想MHD的$\E+\vct u\times\B=0$是其中更强的特殊情形。

所以非零$\E+\vct u\times\B$不一定破坏磁通冻结：它可能只是一个势梯度。反过来，一个局部定义的磁通速度也未必能无奇点地延伸到整个三维重联区域。不能只凭单点非理想电场大小判断全球磁连接关系。

\section{Hall区：磁场可以离开离子而仍随电子运动}
无电阻、无电子惯性的标量电子力平衡是
\begin{equation}
\E+\vct u_e\times\B=-\frac{\nabla p_e}{en}.
\end{equation}
若$p_e=p_e(n)$，称为正压（barotropic）关系；这里“正压”意指压力仅依赖密度，不是说压力数值为正。定义电子压力势
\begin{equation}
h_e(n)=\frac1e\int^n\frac1{n'}\frac{\dd p_e(n')}{\dd n'}\dd n'.
\end{equation}
链式法则给
\begin{equation}
\nabla h_e=\frac1{en}\frac{\dd p_e}{\dd n}\nabla n
=\frac{\nabla p_e}{en}.
\end{equation}
故$\E+\vct u_e\times\B=-\nabla h_e$满足电子磁冻结条件。此时以离子速度写$\E+\vct u_i\times\B$却含Hall项，通常不为纯势梯度。因此“磁场不随离子冻结”与“磁场不随任何流体冻结”是两件不同的事。

\subsection{等温电子的明确例子}
若$\Theta_e$空间恒定，$p_e=n\Theta_e$，则
\begin{equation}
\frac{\nabla p_e}{en}=
\frac{\Theta_e}{e}\nabla\ln n
=\nabla\left[\frac{\Theta_e}{e}\ln\frac n{n_*}\right].
\end{equation}
$n_*$只是使对数无量纲的任意参考密度。这一压力电场可以很强，却不通过其旋度破坏电子磁通冻结。

\subsection{非正压电子与Biermann源}
若$p_e$并非仅为$n$的函数，利用$\nabla\times(f\nabla g)=\nabla f\times\nabla g$可得
\begin{align}
\nabla\times\left(\frac{\nabla p_e}{en}\right)
&=\nabla\left(\frac1{en}\right)\times\nabla p_e\\
&=-\frac{\nabla n\times\nabla p_e}{en^2}
=-\frac{\nabla n\times\nabla\Theta_e}{en}.
\label{hall:eq:biermann}
\end{align}
最后一步用了$p_e=n\Theta_e$。只有密度与温度梯度不平行时才产生此源。若二维梯度全在$x,y$平面，叉积沿$z$，这一项直接改变面外磁场；不要因此误称它已经支撑改变面内连接关系所需的$E_z$。

\section{二维X点的重联电场由哪些项支撑}
采用本讲义坐标：$x$沿出流，$y$沿入流，$z$沿电流或导场，$\partial_z=0$。面内磁通函数取
\begin{equation}
\B_\perp=\nabla A_z\times\hat{\vct z}
=(\partial_yA_z,-\partial_xA_z,0).
\end{equation}
在不含$z$依赖的规范下，$E_z=-\partial_tA_z$。磁岛的磁通变化使用$A_z(O)-A_z(X)$之类的差值，因而还须区分O点与X点的电场；任意单点$E_z$未必就是全局磁通转移率。

对面内X点，$B_x=B_y=0$，而
\begin{equation}
(\vct u_e\times\B)_z=u_{ex}B_y-u_{ey}B_x=0.
\end{equation}
即使有强导场$B_z$，也不会改变这一代数事实。同样，$(\J\times\B)_z=J_xB_y-J_yB_x=0$。故电子欧姆定律在X点的$z$分量为
\begin{align}
E_z(X)={}&-\frac{\partial_xP_{e,xz}+\partial_yP_{e,yz}}{en}\nonumber\\
&-\frac{m_e}{e}\left(\partial_tu_{ez}
+u_{ex}\partial_xu_{ez}+u_{ey}\partial_yu_{ez}\right)
+\eta J_z.
\label{hall:eq:xpoint-ez}
\end{align}
若压力标量，$P_{e,xz}=P_{e,yz}=0$，二维又使$\partial_zp_e=0$。如果同时删去电子惯性、电阻及其他非理想项，则$E_z(X)=0$。这说明Hall项能改变离子尺度结构，却不能单独提供这个二维X点的破冻结电场。\citep{Wang2015}

\subsection{压力值为零不等于压力梯度为零}
在X点附近给一个局部教学构造：$P_{e,xz}=ax$、$P_{e,yz}=by$，取足够大的正对角压力使局部张量正定。在原点两项非对角压力本身都等于零，但
\begin{equation}
E_z(0)=-\frac{a+b}{en}
\end{equation}
可以非零。因此诊断时需要空间导数，不是只看X点压力值。这个构造仅演示局部支撑项，不是一份完整、自洽的重联解。

若X点又是对称电子面内停滞点，$u_{ex}=u_{ey}=0$，那么对流惯性也在该点消失。标量压力、零电阻时剩下$E_z=-(m_e/e)\partial_tu_{ez}$。严格稳态再使时间导数为零，便没有有限$E_z$。非稳态惯性重联不受这一稳态论证排除；不能在推导中先删掉时间项，随后又靠同一时间项解释结果。

\section{电子惯性可以破坏普通磁冻结而保存广义磁冻结}
这个小节解释一个容易误解的概念：发生拓扑变化不一定意味着有正定电阻耗散。仍取电子正压、无碰撞，但保留电子惯性。电子欧姆式为
\begin{equation}
\E+\vct u_e\times\B
=-\nabla h_e-\frac{m_e}{e}\partial_t\vct u_e
-\frac{m_e}{e}(\vct u_e\cdot\nabla)\vct u_e.
\end{equation}
定义电子涡量$\vct\omega_e=\nabla\times\vct u_e$，用恒等式
\begin{equation}
(\vct u_e\cdot\nabla)\vct u_e
=\nabla\frac{u_e^2}{2}-\vct u_e\times\vct\omega_e.
\end{equation}
对整个欧姆式取负旋度，所有梯度旋度为零：
\begin{equation}
\partial_t\B
=\nabla\times(\vct u_e\times\B)
+\frac{m_e}{e}\partial_t\vct\omega_e
-\frac{m_e}{e}\nabla\times(\vct u_e\times\vct\omega_e).
\end{equation}
把时间导数移到一起：
\begin{equation}
\boxed{\partial_t\B_e=\nabla\times(\vct u_e\times\B_e),
\qquad\B_e=\B-\frac{m_e}{e}\vct\omega_e.}
\label{hall:eq:generalized-freezing}
\end{equation}
普通磁场$\B$可以改变连接，而这个广义场仍可满足理想输运。电子惯性提供了新的储能与运动自由度，并不自动成为不可逆热源。

如果还取固定均匀离子、常密度、低频Ampère关系，那么$\vct u_e=-\nabla\times\B/(\mu_0en)$，故
\begin{equation}
\vct\omega_e=-\frac1{\mu_0en}\nabla\times\nabla\times\B
=\frac{\nabla^2\B}{\mu_0en},\qquad
\B_e=\B-d_e^2\nabla^2\B.
\end{equation}
这个常见形式比式\eqref{hall:eq:generalized-freezing}多了固定离子和常密度假设。变密度时旋度还作用于$1/n$，不能仅把$d_e$替换成空间函数便照搬常系数公式。

\section{Poynting定理逐步推导：磁能流去哪里}
Faraday定律点乘$\B/\mu_0$：
\begin{equation}
\partial_t\frac{B^2}{2\mu_0}
=-\frac{\B}{\mu_0}\cdot(\nabla\times\E).
\end{equation}
Ampère--Maxwell方程点乘$\E$：
\begin{equation}
\partial_t\frac{\epsilon_0E^2}{2}
=\frac{\E}{\mu_0}\cdot(\nabla\times\B)-\J\cdot\E.
\end{equation}
将二式相加，再使用
\begin{equation}
\nabla\cdot(\E\times\B)
=\B\cdot(\nabla\times\E)-\E\cdot(\nabla\times\B),
\end{equation}
得到
\begin{equation}
\boxed{\partial_t\left(\frac{\epsilon_0E^2}{2}+\frac{B^2}{2\mu_0}\right)
+\nabla\cdot\left(\frac{\E\times\B}{\mu_0}\right)
=-\J\cdot\E.}
\label{hall:eq:poynting}
\end{equation}
$\vct S=\E\times\B/\mu_0$是Poynting能流，单位W/m$^2$。$\J\cdot\E>0$表示局地电磁场向粒子总能供能；它既可成为流能，也可随后经压力功进入内能，因此不等于全部变成热。

在低频模型中，不推进电场能$\epsilon_0E^2/2$；但仍可由Faraday与无位移Ampère得到磁能形式
$\partial_t(B^2/2\mu_0)+\nabla\cdot\vct S=-\J\cdot\E$。这是与该约化阶次一致的场能账本。

\section{把两物种能量与场能相加}
前面取矩已得到$\mathcal E_s=K_s+U_s$的粒子能量方程。两种物种电功相加恰为
\begin{equation}
en_i\vct u_i\cdot\E-en_e\vct u_e\cdot\E=\J\cdot\E.
\end{equation}
因此与式\eqref{hall:eq:poynting}相加后，这项抵消：
\begin{align}
\partial_t\mathcal E_{\rm all}+\nabla\cdot\vct F_{\rm all}
&=\sum_{s=i,e}(\vct u_s\cdot\vct R_s+H_s),\\
\mathcal E_{\rm all}&=\mathcal E_i+\mathcal E_e+\frac{\epsilon_0E^2}{2}+\frac{B^2}{2\mu_0},\\
\vct F_{\rm all}&=\sum_{s=i,e}(\mathcal E_s\vct u_s+\mathsf P_s\cdot\vct u_s+\vct q_s)
+\frac{\E\times\B}{\mu_0}.
\label{hall:eq:total-energy}
\end{align}
内部保守碰撞必须使右边为零。开放边界下，总能量可改变，但改变应由边界通量解释；外部驱动、辐射、粒子源要单列，不能把一切总能下降都当作“耗散”。

\subsection{摩擦怎样从相对流动取能}
令$\vct R_e=\alpha(\vct u_i-\vct u_e)$、$\vct R_i=-\vct R_e$，其中$\alpha>0$。则
\begin{align}
\vct u_i\cdot\vct R_i+\vct u_e\cdot\vct R_e
&=-(\vct u_i-\vct u_e)\cdot\vct R_e\\
&=-\alpha|\vct u_i-\vct u_e|^2=-\eta J^2.
\end{align}
孤立系统必须有$H_i+H_e=\eta J^2$回到内能。只添加电阻磁扩散而不补热，可能是一个有意设置的冷却模型，但不能称作完整的总能量守恒模型。

\section{Hall项不直接耗散，但可以改变能量路径}
由向量叉积垂直于原向量，
\begin{equation}
\J\cdot\E_H=\J\cdot\frac{\J\times\B}{en}=0.
\end{equation}
这只说明Hall欧姆项对总电磁功的贡献为零。对物种分别看，
\begin{align}
W_{i,H}&=en\vct u_i\cdot\E_H,\\
W_{e,H}&=-en\vct u_e\cdot\E_H,\\
W_{i,H}+W_{e,H}&=en(\vct u_i-\vct u_e)\cdot\E_H=0.
\end{align}
前两项可分别非零、互相抵消。Hall电场结构还改变流路、压缩、焓流与Poynting流的空间分布，从而间接改变热在哪里产生与被带走。

举一个纯代数例：取$\J=J\hat{\vct z}$、$\B=B\hat{\vct x}$，则$\E_H=(JB/en)\hat{\vct y}$。若两种流体有同一个非零$y$速度$U_y$，其Hall电功分别为$+JU_yB$和$-JU_yB$，而总和为零。这里共同$y$速度与沿$z$的电流漂移并不冲突。这个局部例子说明物种交换的可能性，不是完整重联解。

\section{双温Hall模型的总能量通量还包含电子焓流}
这一步可检查上一章的两温方程是否自洽。无电阻、无热流时，记$K=m_inu^2/2$、$U_i=3p_i/2$、$U_e=3p_e/2$。由离子尺度总动量和两套内能相加，可整理为
\begin{align}
&\partial_t(K+U_i+U_e)
+\nabla\cdot[(K+U_i+p_i)\vct u+(U_e+p_e)\vct u_e]\nonumber\\
&\hspace{20mm}=\vct u\cdot(\J\times\B)-\frac{\J}{en}\cdot\nabla p_e
=\J\cdot\E.
\label{hall:eq:hall-particle-energy}
\end{align}
为看清电子焓流怎样出现，先把压力总梯度写成$-\nabla\cdot[(p_i+p_e)\vct u]+(p_i+p_e)\nabla\cdot\vct u$。与两物种压缩功相加后，剩下$p_e\nabla\cdot(\vct u-\vct u_e)$。再用
\begin{equation}
p_e\nabla\cdot\vct w=\nabla\cdot(p_e\vct w)-\vct w\cdot\nabla p_e,
\qquad\vct w=\frac{\J}{en},
\end{equation}
将第一个散度移入通量，原来的电子压力通量$p_e\vct u$就变为$p_e\vct u_e$，得到式\eqref{hall:eq:hall-particle-energy}。

再与磁能式相加，总能密度$K+U_i+U_e+B^2/(2\mu_0)$守恒。若把电子内能随离子速度搬运却不加相应修正，便改变了这个通量结构。等温电子模型另有热库；它不应直接套用本节的孤立绝热总能公式。

\section{压缩、剪切与不可逆热化必须分别判断}
将压力写成$\mathsf P=p\mathsf I+\boldsymbol{\Pi}$，其中$\operatorname{tr}\boldsymbol{\Pi}=0$。再将速度梯度分为各向同性膨胀、无迹对称应变$\mathsf D$和反对称旋转。压力对反对称部分的缩并为零，所以
\begin{equation}
-\mathsf P:\nabla\vct u=-p\theta-\boldsymbol{\Pi}:\mathsf D,
\qquad\theta=\nabla\cdot\vct u,
\quad D_{ij}=\frac{\partial_ju_i+\partial_iu_j}{2}-\frac13\theta\delta_{ij}.
\end{equation}
第一项是压缩功，第二项是偏离各向同性压力与应变的能量交换。无碰撞时二者都不必逐点正定；内能增大也可能是可逆变形。

对于真正的Newton黏性闭合$\boldsymbol{\Pi}=-2\mu\mathsf D$且$\mu>0$，第二项才成为$2\mu\mathsf D:\mathsf D\geq0$。类似地，正电阻给$\eta J^2\geq0$。这些符号来自特定本构关系，不是由“压力功”这个名称保证。

对光滑、标量、绝热物种，令每粒子无量纲熵为
\begin{equation}
\sigma_s=\frac32\ln\Theta_s-\ln n_s+\text{常数}
=\frac32\ln\left(\frac{p_s}{n_s^{5/3}}\right)+\text{常数},
\end{equation}
对数中的量理解为除以固定参考量。连续性和压力方程给$D_s\sigma_s=0$，即使压缩使温度显著升高。这是“增温不必增熵”的明确演算。碰撞少的多峰分布，其精细分布熵不是$n,\vct u,p$几个量就能唯一恢复；有限矩熵和动理学熵也应区分。

\note{报告能量分配时至少分清：物种电功、平均流能、内能、局部温升、焓流、非热尾与不可逆熵产。不同定义得到不同的离子/电子比并不自动矛盾；必须在同一控制体、同一时间窗、同一种物理量上比较。}
